\documentclass[12pt]{article}

\usepackage[utf8]{inputenc}
\usepackage{amsmath, amssymb, amsthm}
\usepackage{geometry}
\geometry{a4paper, margin=2.5cm}

\title{Teorema Casorati--Weierstrass}
\author{}
\date{}

\theoremstyle{definition}
\newtheorem{definition}{Definiție}

\theoremstyle{plain}
\newtheorem{theorem}{Teoremă}
\newtheorem{proposition}{Propoziție}

\begin{document}

\maketitle

\begin{definition}
Fie $f$ o funcție olomorfă în domeniul perforat


\[
0 < |z - z_0| < r.
\]


Spunem că $z_0$ este o \emph{singularitate esențială} a lui $f$ dacă seria Laurent a lui $f$ în jurul lui $z_0$ conține o infinitate de termeni negativi, sau echivalent, dacă singularitatea nu este nici eliminabilă, nici pol.
\end{definition}

\begin{theorem}[Casorati--Weierstrass]
Fie $z_0$ o singularitate esențială a funcției olomorfe $f$. Atunci imaginea lui $f$ în orice vecinătate perforată a lui $z_0$ este densă în $\mathbb{C}$.

Mai precis, pentru orice $w \in \mathbb{C}$ și orice $\varepsilon > 0$, există $z$ cu $0 < |z - z_0| < \varepsilon$ astfel încât


\[
f(z) \in B(w, \varepsilon).
\]


\end{theorem}

\begin{proposition}
În ipotezele teoremei, următoarele afirmații sunt echivalente:
\begin{enumerate}
    \item[(i)] $z_0$ este o singularitate esențială a lui $f$;
    \item[(ii)] $f$ nu este mărginită în nicio vecinătate perforată a lui $z_0$;
    \item[(iii)] Imaginea lui $f$ în orice vecinătate perforată a lui $z_0$ este densă în $\mathbb{C}$.
\end{enumerate}
\end{proposition}

\end{document}
